\section{The two supports of the arithmetic theta column}
\label{csp-section}

We identify the finite operator defining the column with an
ordinary measure on each of the two required matrix supports.
The calculation keeps the normalized Fourier denominator,
the positive square root in the lattice move, and the
determinant character.  Their combination gives the precise
factor $p^{-fa}$ in the selected type.

\subsection{The selected prime and the actual moment coordinates}

Use the fields, ideal $I$ of norm $A$, inverse half different
$\mathfrak a$, and normalized family of
Section~\ref{te-section}; here the capital $A$ distinguishes
the elliptic ideal norm from the matrix rank $a$.
Let $a,b$ be positive even integers and let
$N\subset E^a$, $N'\subset E^b$ be unimodular Hermitian
$\mathcal O_E$-lattices.  Put
\[
 \mathcal M=I\mathfrak a\operatorname{Hom}_{\mathcal O_E}(N,N'),
 \qquad \Lambda=\delta_I\operatorname{emb}(\mathcal M),
                         \qquad F=e_\chi(N,N'),              \tag{CS1}
\]
where $\delta_I$ is the discriminant twelfth root used for
the elliptic ideal lattice.  All embeddings in $\operatorname{emb}$
belong to the CM type $\Sigma$, and $f=|\Sigma|$.

The selected prime data are
\[
 p\ge5,\quad \pi\in\mathcal O_E,\quad
 \pi\bar\pi=p,\quad (\pi,\bar\pi)=1,\quad
 \mathcal P=\{P:P\mid\pi\},\quad
                 |\mathcal P|=f,\quad N(P)=p.              \tag{CS2}
\]
Assume $p\nmid A$ and that $I\mathfrak a$ is a unit at
every prime over $p$.  These are the actual selected-prime
conditions.  In particular $p$ splits completely in $E$:
$N_{E/\mathbb Q}(\pi)=p^f$, so each of the $f$ primes in
$\mathcal P$ has multiplicity one; the conjugate primes
are disjoint and account for the other $f$ primes over $p$.
The inverse half different is a unit there because its inverse
is an integral ideal whose product with its conjugate is the
relative different, and this different is a unit at a completely
split prime.  Thus (CS1) is locally the ordinary Hom lattice
at these primes.

For each $P\in\mathcal P$ fix nonzero residue functionals
\[
 \ell_P:N_P\longrightarrow\mathbb F_p,\qquad
              \ell'_P:N'_P\longrightarrow\mathbb F_p.       \tag{CS3}
\]
They factor through reduction modulo $P$ and are surjective.
They define the finite local test function
\[
 \phi_P(X)=\mathbf1_{\{\ell'_PX=\ell_P\}},\qquad
             X\in\operatorname{Hom}_{\mathbb F_p}
                       (N_P/PN_P,N'_P/PN'_P).               \tag{CS4}
\]
Outside the integral Hom lattice its value is zero.
Let $\phi=\prod_{P\in\mathcal P}\phi_P$.

There are two types of transported place used below.
In the selected type, every map $\iota\eta$, $\eta\in\Sigma$,
induces a prime in $\mathcal P$.  The ordinary measure
variable in (GM8) then consists of the conjugate local
components at $\overline{\mathcal P}$.
In the conjugate type, every $\iota\eta$ induces a prime
in $\overline{\mathcal P}$ and the measure variable consists
of the components at $\mathcal P$.  These are precisely the
two place conditions used for the two column bounds.

We record the coordinate constants explicitly.  Let $B$ be a
number field containing all embeddings of $E$ in $\Sigma$,
and put
\[
 \mathcal L_B=\mathcal O_B\langle(\eta(x))_{\eta\in\Sigma}:
                                      x\in\mathfrak a\rangle,
 \qquad J_\eta=\{c\in B:c\,e_\eta\in\mathcal L_B\}.
                                                               \tag{CS5}
\]
The prescribed eigen-direction generator $c_\eta$ satisfies
$J_\eta=c_\eta\mathcal O_B$, after the finite extension in
the arithmetic datum.

\begin{lemma}\label{csp-coordinate-units}
At every transported place above $p$ one has
$|\iota(c_\eta)|_p=1$.
With the common period $\omega=\Omega_0/\iota(\delta_0)$,
put
\[
 \Omega_I=A\iota(\delta_I)\omega
       =A\Omega_0\iota(\delta_I/\delta_0),\qquad
            u_\eta=\frac{\iota(c_\eta)}{\Omega_I}.           \tag{CS6}
\]
All $u_\eta$ have absolute value one.  The moment of the
mixed derivative with direction $c_\eta$ in entry $(i,j)$
is exactly the corresponding product of the coordinates
\[
                         u_\eta\,\iota\eta((X^*)_{ji}).     \tag{CS7}
\]
There are no factorials in this mixed-derivative convention.
\end{lemma}
\begin{proof}
Every coordinate of $\mathcal L_B$ is integral at the chosen
place, since $\mathfrak a$ is a unit above $p$.  Hence
$|\iota(c_\eta)|_p\le1$.
For $\eta'\ne\eta$, complete splitting and the Chinese
remainder theorem give $x\in\mathcal O_E$ for which
$\iota\eta(x)-\iota\eta'(x)$ is a unit.  Start with the
vector $(1)_\eta\in\mathcal L_B$, which belongs to this
lattice because $1\in\mathfrak a$.  For each $\eta'\ne\eta$
apply the operator
\[
         (v_\rho)_\rho\longmapsto
                     ((\rho(x)-\eta'(x))v_\rho)_\rho.       \tag{CS8}
\]
It preserves $\mathcal L_B$, kills the $\eta'$ coordinate,
and multiplies the $\eta$ coordinate by a unit at the chosen
place.  After all the operators, the resulting vector is
$c e_\eta$ with $c\in J_\eta$ and $|\iota(c)|_p=1$.
Writing $c=c_\eta v$ with $v\in\mathcal O_B$ gives the
reverse inequality for $|\iota(c_\eta)|_p$.

The period and discriminant ratio in (CS6) are units by
Proposition~\ref{ci-common-period} and the ideal discriminant
comparison in that section; $p\nmid A$.  This proves the
unit assertions.
For a coordinate direction $v=c_\eta e_{\eta,i,j}$,
the Hermitian form in (TE2) gives, on rational lattice points,
\[
 h\bigl(v,\delta_I\operatorname{emb}(X)\bigr)
                 =\frac{c_\eta}{A\delta_I}\,
                                      \eta((X^*)_{ji}).      \tag{CS9}
\]
Multiply by $\omega^{-1}$ as in the derivative form of
(GM13).  This gives (CS7), with (CS6) exactly as stated.
Multilinearity proves the mixed-derivative formula.  These
coordinate directions lie in the algebraic linear span of
$\delta_0^{-1}\Lambda$: the embedding matrix is algebraic
and invertible, and $\delta_I/\delta_0$ is algebraic.
Corollary~\ref{gm-algebraic-jets} therefore also proves
algebraicity of these mixed derivatives, including all the
rational translates occurring below.
\end{proof}

\subsection{The finite operator, with its normalization}

Define the two commensurable self-dual lattices
$\mathcal A,\mathcal B\subset\mathcal M\otimes\mathbb Q$
by their localizations.  At a selected prime $P$, put
\[
 \mathcal A_P=p^{-1}\mathcal M_P,\quad
 \mathcal B_P=p\mathcal M_P,\qquad
 \mathcal A_{\bar P}=p\mathcal M_{\bar P},\quad
 \mathcal B_{\bar P}=p^{-1}\mathcal M_{\bar P};              \tag{CS10}
\]
at other primes both equal $\mathcal M$.
In formulas with $U,W,E$, identify rational matrices with
their scaled embedded vectors in (CS1).
The finite groups
\[
 T=(\mathcal M\cap\mathcal A)/(\mathcal A\cap\mathcal B),
 \qquad Y=\mathcal B/(\mathcal M\cap\mathcal B)             \tag{CS11}
\]
have order $p^{fab}$ and pair perfectly by
$\kappa(t,y)=e^{-2\pi iE(t,y)}$.
The first is the product of the matrix spaces in (CS4).
This follows directly from (CS10); the local trace duality
is perfect because $\mathcal M$ is self-dual.

For a function $\Phi$ on $T$, the finite theta operator is
\[
 \mathcal Q_\Phi F=
       \frac1{|T|}\sum_{y\in Y}\sum_{t\in T}
                  \Phi(t)\kappa(t,y)^{-1}W_yF.              \tag{CS12}
\]
Representatives in the indicated lattices are understood.
The expression is independent of them: on $\mathcal B$ the
operators $W_y$ are additive; representatives differing by
$\mathcal M\cap\mathcal B$ act identically on $F$; and
the characters factor through (CS11).

Here is the precise connection with the actual theta indicator.
At a general level $c$ and selected set of primes, use the
two half lattices defined as in (CS10), with $c$ in place
of $p$.  Put
$n=[\mathcal M\cap\mathcal A:\mathcal A\cap\mathcal B]$
and take representatives $R$ of
$\mathcal B/(\mathcal M\cap\mathcal B)$.
For a rational selected coset $x$, choose a lift $s_x$ in
some half lattice $\mathcal A(c')$, with $c'>0$, whose
selected components agree with $x$ modulo $c\mathcal M$.
The Chinese remainder theorem supplies such a lift.
The actual indicator is
\[
                  \Theta_x=\frac1n W_{-s_x}\sum_{y\in R}W_yF.
                                                               \tag{CS12a}
\]
Indeed its definition has an outside factor
$\mathrm{vol}^{-1/2}$.  Its Fourier coefficient at the
indicator of $\mathcal B$ is
$\mathrm{vol}^{-1/2}(\mathrm{vol}/n)$: only the zero
class of $\mathcal M\cap\mathcal A$ contributes.  The
two factors cancel exactly, giving $1/n$ in (CS12a).
Here $\mathrm{vol}>0$, since all lattice indices are
positive.  No assertion that this volume is a $p$-adic
unit is involved.

The lift in (CS12a) is independent of choice.  Two lifts
$s,t$ differ by an element of $\mathcal B$.  Moreover they
both belong to a common integral half lattice: if their
levels are $c',c''$, use $\mathcal A(\gcd(c',c''))$.
At a selected prime their difference is integral, so the
smaller of the two possible denominator bounds bounds
both; at the conjugate prime both have the corresponding
integral divisibility.  At other primes they belong to
$\mathcal M$.  Consequently $W_{-t}W_{t-s}=W_{-s}$
by the semicharacter law on that half lattice.  The sum
over $R$ is fixed by $W_{t-s}$, by permutation of its
classes.  This proves the claimed independence, including
its phase.

At $p$, the cosets with nonzero (CS4) are the elements
of $T$ in (CS11).  Choose their lifts in $\mathcal M$.
Commuting $W_{-s_x}$ through the sum in (CS12a), and using
$W_{s_x}F=F$, gives exactly (CS12).  Thus its denominator
is $|T|$, with no further $p$-primary factor.

\begin{lemma}\label{csp-fourier-fiber}
Choose $X_0\in\mathcal M$ satisfying
$\ell'_PX_0=\ell_P$ modulo $P$ for every selected $P$,
and put $K=\ker(\prod_P(\ell'_P\cdot))\subset T$.
Let $K^\perp\subset Y$ be its annihilator.  Then
\[
 |T/K|=|K^\perp|=p^{fa}=:h_0,
 \qquad
 \frac1{|T|}\sum_{t\in T}\phi(t)\kappa(t,y)^{-1}
       =\begin{cases}
          h_0^{-1}\kappa(X_0,y)^{-1},&y\in K^\perp,\\
          0,&y\notin K^\perp.
        \end{cases}                                       \tag{CS13}
\]
\end{lemma}
\begin{proof}
For each $P$, choose $v\in N'_P/PN'_P$ with
$\ell'_P(v)=1$.  The map $X=v\ell_P$ is a solution of
the affine equation.  Chinese remaindering lifts all these
solutions to one $X_0\in\mathcal M$.
The linear map $X\mapsto\ell'_PX$ onto the $a$-dimensional
dual of $N_P/PN_P$ is surjective, so its kernel has
codimension $a$.  Hence $|T/K|=p^{fa}$.
The support of $\phi$ is the coset $X_0+K$.
Summing a character on this coset gives zero unless it is
trivial on $K$, in which case it gives
$|K|\kappa(X_0,y)^{-1}$.  Division by $|T|$ proves (CS13).
The equality $|K^\perp|=|T/K|$ follows from perfectness
of the pairing in (CS11).
\end{proof}

\subsection{The conjugate type: exact restriction of the measure}

Suppose first that $\iota$ has the conjugate type.  Then
$Y$ is the formal $p$-torsion component, and the measure
variable is the selected component of $\mathcal M$.
Representatives of $Y$ can be chosen among the strong good
lifts in (GM10), with any prescribed additional
prime-to-$p$ multiple in $S$.  To see this explicitly,
choose $\kappa\in\mathcal O_k$ whose formal component is
$S^{-1}$ modulo $p^2$ and whose conjugate component is zero
modulo $p^2$.  The vectors $S\kappa l/p$, $l\in\Lambda$,
have selected component in $p\mathcal M_P$, lie in
$\mathcal B$, and realize every class of $Y$.
They differ from any other representatives by
$\mathcal M\cap\mathcal B$, so the operator is unchanged.
The exact law (GM14) therefore identifies (CS12) with the
measure
\[
 \left(\frac1{|T|}\sum_{y\in Y}\sum_{t\in T}
          \phi(t)\kappa(t,y)^{-1}\kappa(X,y)\right)d\mu_F(X)
                              =\phi(X)\,d\mu_F(X).          \tag{CS14}
\]
The equality is finite Fourier inversion.  In particular its
norm is at most one, with support
\[
 S_{\bar P}=\{X:\ XN_P\subset N'_P,\quad
                        \ell'_PX=\ell_P\text{ for all }P\in\mathcal P\}.
                                                               \tag{CS15}
\]
The Fourier denominator has disappeared by exact character
orthogonality.  Bounding its terms separately would give an
incorrect and unnecessarily large denominator.

\subsection{The selected type: the companion lattice}

Define a new unimodular lattice $N''$ by
\[
 N''_P=\ker\ell'_P\subset N'_P,\qquad
 N''_{\bar P}=(N''_P)^\#,\qquad
 N''_Q=N'_Q\quad(Q\nmid p).                               \tag{CS16}
\]
The dual at $\bar P$ is for the perfect Hermitian pairing
between the two conjugate completions.
Local patching, as in Lemma~\ref{acl-patching}, gives a
global projective lattice with these completions; their
duality proves global unimodularity.  Put
\[
           \mathcal M''=I\mathfrak a
                         \operatorname{Hom}_{\mathcal O_E}(N,N'').
                                                               \tag{CS17}
\]

\begin{lemma}\label{csp-companion}
The group $K^\perp$ in (CS13) is the image of
$\mathcal B\cap\mathcal M''$ in $Y$, and this image
identifies with
$\mathcal M''/(\mathcal M\cap\mathcal M'')$.
This quotient has order $h_0=p^{fa}$.  Moreover
\[
        \det(N',N'')=(\pi/\bar\pi),\qquad
        \pi N'\subset N'',\quad \bar\pi N''\subset N'.     \tag{CS18}
\]
\end{lemma}
\begin{proof}
At $P$, the kernel of $X\mapsto\ell'_PX$ consists of maps
whose image is in $\ker\ell'_P$.  Under the perfect matrix
trace pairing, its annihilator consists of adjoint matrices
that take $\ker\ell'_P$ integrally into $N_P$.
Equivalently, in bases with $\ell'_P$ the last coordinate,
the only possibly nonintegral row of the adjoint matrix is
its last row, with denominator $p$.  The other rows are
integral.  This is exactly the dual of
$\operatorname{Hom}(N_P,\ker\ell'_P)$ at $\bar P$,
which is $\mathcal M''_{\bar P}$.
It enlarges $\mathcal M_{\bar P}$ by an index $p^a$.

At $P$ the lattice $\mathcal M''_P$ is contained in
$\mathcal M_P$ and contains $p\mathcal M_P$.
Thus its selected component creates no quotient modulo
$\mathcal M$.  By the Chinese remainder theorem one may
choose representatives of $\mathcal M''/(\mathcal M\cap\mathcal M'')$
whose selected components are in $p\mathcal M_P$.
They then belong to $\mathcal B$ and realize the indicated
annihilator.  The quotient size is $\prod_Pp^a=h_0$.

The relative determinant has valuation $1$ at each $P$,
$-1$ at each $\bar P$, and zero elsewhere.  These are
the valuations of $(\pi/\bar\pi)$ by (CS2), proving the
first equality in (CS18).  The two inclusions follow
locally from (CS16): at $P$ the kernel contains $pN'_P$,
and at $\bar P$ its dual is contained in $p^{-1}N'_{\bar P}$.
At all other primes the lattices agree and $\pi,\bar\pi$
are units.
\end{proof}

The exact right-move identity (TE17) now gives
\[
 \sum_{y\in K^\perp}W_yF
      =\sqrt{h_0}\,\chi((\pi/\bar\pi))^{a/2}
                                      e_\chi(N,N'').        \tag{CS19}
\]
The positive real square root is the same one as in the
definition of the averaging operator.  Since $X_0\in\mathcal M$,
$W_{X_0}F=F$, and the Heisenberg commutator yields
\[
 \mathcal Q_\phi F
    =h_0^{-1}W_{-X_0}\sum_{y\in K^\perp}W_yF
    =h_0^{-1}\sqrt{h_0}\,
       \chi((\pi/\bar\pi))^{a/2}
                              W_{-X_0}e_\chi(N,N'').       \tag{CS20}
\]
This proves the needed identity of theta functions, before
taking a valuation.

Suppose $\iota$ has the selected type.  The character convention
(TE16) and (CS2) give
\[
 \left|\iota\chi((\pi/\bar\pi))\right|_p=p^f,\qquad
 \left|\iota\left(h_0^{-1}\sqrt{h_0}
                         \chi((\pi/\bar\pi))^{a/2}\right)\right|_p
                                                        =p^{fa}.
                                                               \tag{CS21}
\]
Indeed every $\iota\eta(\pi)$ has absolute value $p^{-1}$,
and every $\iota\eta(\bar\pi)$ has absolute value one.
The square root has absolute value $p^{-fa/2}$, cancelling
the absolute value $p^{fa/2}$ of the character factor.

The class of $X_0$ modulo $\mathcal M''$ is formal
$p$-torsion in this type: its selected component may have
order $p$, while its conjugate component already lies in
$\mathcal M''_{\bar P}$.  We spell out the harmless
choice of its representative.  If $r$ is any rational
formal $p^n$-torsion lift modulo a self-dual lattice, choose
a strong good lift $y$ congruent to $r$ modulo that lattice,
using the two ordinary idempotents and the Chinese remainder
theorem as above.  For $l=r-y$ in the lattice,
the Heisenberg formula and the lattice theta multiplier give
\[
                         W_rF'=\zeta W_yF',                 \tag{CS22}
\]
where $\zeta$ is a root of unity, since all quadratic and
alternating values at rational lattice points are rational.
The exact law for $y$ therefore gives a measure of norm
at most one for $W_rF'$, on the same ordinary lattice
component as for $F'$.  This observation needs only the
absolute value of $\zeta$, not an unspecified equality of
phases in (CS20).

Apply this to $r=-X_0$ and $F'=e_\chi(N,N'')$.
By (CS20)--(CS22), the selected-type measure for
$\mathcal Q_\phi F$ has norm at most $p^{fa}$ and is
supported on the conjugate ordinary component of
$\mathcal M''$.  In the original matrix language its
support is contained in
\[
 S_P=\{X:\ X^*(\ker\ell'_P)\subset N_P
                               \text{ for all }P\in\mathcal P\}.
                                                               \tag{CS23}
\]
To check the inclusion, let $z\in\ker\ell'_P=N''_P$.
For $v\in N_{\bar P}$, the image $Xv$ lies in $N''_{\bar P}$
on the conjugate component.  Perfect Hermitian duality
gives $h(v,X^*z)=h(Xv,z)\in\mathcal O_{E,P}$.
Since $N_P$ is its appropriate local dual, $X^*z\in N_P$.
This is exactly the condition in (CS23).

\subsection{The other primes and simultaneous moment approximation}

We verify that restoring the other local test functions does
not change either bound.  Let $S_0$ be the finite set of
auxiliary primes, disjoint from the primes over $p$, and
assume $S_0\cap\overline{S_0}=\varnothing$, as in the
selected local datum.  Let $c_A,g$ be positive integers
prime to $p$.
The full test function has the form
\[
                              \Phi=\phi\,\Phi_0,             \tag{CS24}
\]
where $\Phi_0$ is supported in the prescribed
$c_A^{-1}$ integral Hom lattice at $S_0$, is invariant under
$pg$ times that local Hom lattice, and all its transported
values have absolute value at most one.  Invariance means
that it is a function on the finite quotient at $S_0$;
it does not impose an additional condition at $p$.
These are the away-from-$p$ datum conditions in the column.

\begin{lemma}\label{csp-other-primes}
The full finite theta operator associated to (CS24), with
the same parameter $\mathrm{vol}=1$, is a finite sum
\[
                           \sum_j d_jW_{r_j}\mathcal Q_\phi F,
                                                               \tag{CS25}
\]
where $|\iota(d_j)|_p\le1$, each $r_j$ has order prime to
$p$ modulo both $\mathcal M$ and $\mathcal M''$, and the
representatives may be chosen arbitrarily divisible by $p$
at all primes over $p$.  In the conjugate type its measure
has norm at most one and support (CS15); in the selected
type it has norm at most $p^{fa}$ and support (CS23).
\end{lemma}
\begin{proof}
All representative quotients in the finite formula split
as the direct product of their $p$-primary part and their
part of order prime to $p$.  This is the elementary primary
decomposition of a finite abelian group, or equivalently
the Chinese remainder theorem on each localized lattice.
The alternating pairings between the factors are trivial:
a character pairing value has order dividing the orders
of both factors.  Each normalized Fourier denominator
therefore factors as a power of $p$ times an integer prime
to $p$.

Use the indicator formula (CS12a), and sum it with the
values of (CS24) over the finite selected coset quotient.
Its $p$-part is $T$; the $p$-part of the representative
group $R$ is $Y$; and the $p$-part of $n$ is $|T|$.
These assertions follow directly from (CS10), since
$p\nmid c_Ag$.  Choose $s_x=s_{x_p}+s_{x_0}$ according
to the primary decomposition, with $s_{x_p}\in\mathcal M$
and $s_{x_0}$ arbitrarily $p$-divisible.  Such representatives
exist by Chinese remaindering.  Expand the remaining
prime-to-$p$ finite sums in (CS12a).  Each coefficient is
a value of $\Phi_0$, times roots of unity and the inverse
of an integer prime to $p$, so has absolute value at most
one.  The remaining $p$-primary sum is exactly (CS12),
because (CS24) factors and its selected part is (CS4).

Here are explicit compatible choices of representatives.
If an away-from-$p$ class has order dividing $n$, choose
an integer $j\equiv1\pmod n$ with $j\equiv0\pmod{p^m}$,
and replace its representative $r$ by $jr$.
The class is unchanged, and its $p$-components are as
divisible as desired.  Choose the $p$-primary representatives
with the additional factor $S$ in (GM10) divisible by all
the finitely many away-from-$p$ denominators.  Their
pairings with the other representatives are then integral
at every prime.  Thus their $W$ operators commute exactly.
Changing a representative in its own integral lattice
quotient does not change the corresponding sum.  Combining
the remaining two operators introduces at most a root of
unity by the rational Heisenberg formula.  This proves
(CS25).  Taking $m\ge2$ ensures in addition that $r_j$
is integral in both lattices $\mathcal M,\mathcal M''$
at $p$, so its order in either quotient is prime to $p$.

For the conjugate type one may instead apply the $p$-primary
Fourier sum after each away-from-$p$ translate.  It then
restricts the norm-one measure of $W_{r_j}F$ by (CS14).
The chosen $r_j$ has zero reduction at the selected primes,
so the same affine condition (CS15) occurs.
For the selected type commute $W_{r_j}$ with $W_{-X_0}$
in (CS20).  The introduced factor is a root of unity.
Apply the ordinary-torsion law to the prime-to-$p$ translate
$W_{r_j}e_\chi(N,N'')$.  One can choose its strong lift
congruent to $-X_0$ modulo an additional prime-to-$p$
multiple of $\mathcal M''$, so the same argument (CS22)
applies.  Its measure still has norm at most one on the
ordinary component of $\mathcal M''$.
Equation (CS21) gives the stated bound and (CS23) the
support.  Finite sums with coefficients bounded by one
preserve these bounds by the ultrametric inequality.
\end{proof}

\begin{theorem}\label{csp-two-supports}
For the actual finite theta column with standard
representatives at the selected primes, its mixed-derivative
table has the simultaneous finite approximation required
in (CV9)--(CV10).  In the conjugate type the approximating
point matrices belong to $S_{\bar P}$ and their coefficients
have absolute value at most one.  In the selected type
they belong to $S_P$ and the table is approximated by
\[
       \sum_\nu b_\nu p^{-fa}
             \prod_{\eta,i,j}
                \bigl(u_\eta\iota\eta((X_\nu^*)_{ji})\bigr)^{m_{\eta,i,j}},
                        \qquad |b_\nu|_p\le1.               \tag{CS26}
\]
The $u_\eta$ are the fixed units (CS6), independent of the
derivative multi-index and the finite approximation.  Thus
the two ordinary support inputs in
Theorem~\ref{cv-two-kernels} are proved for this column.
\end{theorem}
\begin{proof}
Lemmas~\ref{csp-coordinate-units} and
\ref{csp-other-primes} give measures with the exact point
moments and the stated bounds.  Global rational lattice
points are dense in their ordinary completions, by the
Chinese remainder theorem.  In the conjugate type impose
the fixed residue equations (CS15); lifting those equations
and then arbitrarily many additional $p$-adic digits proves
density in the restricted support.  In the selected type
the global points of $\mathcal M''$ are dense in its
ordinary component and all satisfy (CS23).

Apply Theorem~\ref{lm-finite-approximation} to any finite
family of the point polynomials (CS7).  It approximates
all of them at once by a finite sum of point masses,
with coefficients bounded by the measure norm.
For the selected type multiply the measure by $p^{fa}$
first.  Its norm is then at most one, giving exactly
(CS26) after division by $p^{fa}$.

The arithmetic column also divides the theta jets by its
determinant character and inserts the auxiliary local
character values.  For standard representatives these
are $p$-adic units, by the principal-ideal argument
(GM17) and the prime-to-$p$ conductor condition.  The
normalization is therefore a scalar of absolute value
one and preserves both approximation assertions.
Where the coefficient table is descended to its specified
number field, apply the fixed embedding of that field;
algebraicity and the equality of the two descriptions of
the same jet identify the transported values.  This is
the precise descent condition in the actual table and
does not introduce an additional denominator.
\end{proof}

The two bounds above use the actual local affine equation,
the actual companion lattice, and the exact determinant
character.  A bounded measure with no character identities
would not give (CS14), and an eighth-power comparison of
theta vectors would not give (CS19).  Those dependencies
were separately proved in Sections~\ref{gm-section} and
\ref{te-section}.  The incidence identity between adjacent
columns is a further statement and is not asserted here.


